પ્રથમ-ક્રમની ભાષામાં \emph{અચળ}, \emph{ફલન} અને \emph{વિધેય} સંકેતો હોય છે. ફલન અને વિધેય સંકેતો નિર્દિષ્ટ સંખ્યામાં દલીલો લે છે. દાખલા તરીકે, \emph{અંકગણિતની ભાષા}માં એક અચળ સંકેત $\Obj 0$, એક એકસ્થાની ફલન સંકેત $\prime$, બે દ્વિસ્થાની ફલન સંકેતો $+$ અને $\times$ તથા એક દ્વિસ્થાની વિધેય સંકેત $<$ છે. \emph{ચલો}, અચળો અને ફલન સંકેતોથી ભાષાનાં \emph{પદો} બને છે. પદો, વિધેય સંકેતો અને \emph{તાદાત્મ્ય સંકેત} $\eq$ વડે \emph{પરમાણુક સૂત્રો} બનાવીએ છીએ. પછી $ \lnot$, $\lor$, $\land$, $\lif$, $\liff$ જેવા તાર્કિક સંયોજકો અને $\lforall$, $\lexists$ પરિમાણકો વડે તેમાંથી સૂત્રો બનાવીએ છીએ. પદો અને સૂત્રો બનાવતી વખતે જરૂરી કૌંસ હંમેશાં લખીએ છીએ, તેથી સૂત્રને વાંચવાની રીત એક જ હોય છે. આમ સૂત્રની સંરચના પર આગમન વડે વ્યાખ્યાઓ આપી શકાય છે.

સૂત્રમાં ચલની આવૃત્તિ ક્યારેક તેને અનુરૂપ પરિમાણકના નિયંત્રણમાં હોય છે. એ પરિમાણકના \emph{વ્યાપકક્ષેત્ર}માં આવતી આવૃત્તિ \emph{બદ્ધ} છે. બાકી આવૃત્તિઓ \emph{મુક્ત} છે. આ ખ્યાલોની આગમનાત્મક વ્યાખ્યાઓ છે. સૂત્રમાં ચલને સ્થાને પદનું \emph{પ્રતિસ્થાપન} કરવાની ક્રિયા પણ આગમનાત્મક રીતે વ્યાખ્યાયિત કરીએ છીએ. જેમાં ચલની કોઈ મુક્ત આવૃત્તિ ન હોય તેવાં સૂત્રો \emph{વાક્યો} છે.
